%% annex/p12/sat-and-certificates.tex --- the sole new Part-12 annex item

\begin{annexitem}{annex:p12.sat.1}{thm:p12-dpll-correctness}{Correctness of the displayed DPLL logical core}
\begin{proof}
Let the current state be a residual clause set \(G\), a consistent partial
truth assignment \(u\) to variables of the original finite input \(F\), and
the unchanged finite set \(\operatorname{Var}(F)\).  Maintain both parts of the
following invariant: every original clause deleted from \(G\) already contains
a literal made true by \(u\), and every clause retained in \(G\) is exactly the
residual of its original clause after literals made false by \(u\) have been
deleted.

A unit step selects a singleton clause \(\{L\}\).  Every total satisfying
extension of \(u\) must make \(L\) true.  Extending \(u\) accordingly, deleting
clauses containing \(L\), and deleting the complement of \(L\) elsewhere
therefore preserves the invariant and preserves existence of a satisfying
extension in both directions.  At any state, including the initial state, if
\(G\) contains the empty clause, no extension of the current assignment can
satisfy that residual clause; a unit update that exposes it therefore fails.
At any state, including the initial state, if \(G\) has no clauses, each
original clause already has a true literal; an update that deletes the last
residual clause is one instance of this case.  Assigning
arbitrary truth values to the remaining original variables cannot
destroy any such witness, so the returned total assignment satisfies \(F\).

If neither terminal case applies, the selected branch variable \(p\) is not in
the domain of \(u\).  Every total extension makes exactly one of \(p\) and
\(\neg p\) true.  Consequently the current state has a satisfying extension if
and only if at least one of the two child states does.  Apply this equivalence
recursively.  Every unit step and every branch assigns a previously unassigned
member of the fixed finite set \(\operatorname{Var}(F)\), so no branch is
infinite.  Induction on the number of unassigned original variables now gives
both return specifications: a successful leaf returns a total satisfying
assignment, and failure of both children proves that no satisfying assignment
extends the parent.  At the root these are exactly satisfiability and
unsatisfiability of \(F\).
\end{proof}
\end{annexitem}
